\documentclass[11pt,a4paper]{article}
\usepackage[margin=24mm]{geometry}
\usepackage[T1]{fontenc}
\usepackage{lmodern,amsmath,amssymb,amsthm,mathtools,booktabs,longtable,array}
\usepackage[dvipsnames]{xcolor}
\usepackage{microtype,fancyhdr,xurl}
\usepackage[colorlinks=true,linkcolor=MidnightBlue,citecolor=MidnightBlue,urlcolor=MidnightBlue]{hyperref}
\usepackage{enumitem}
\setlist{nosep}
\setlength{\parskip}{5pt}
\setlength{\parindent}{0pt}
\newcommand{\xc}{\operatorname{xc}}
\newcommand{\rankp}{\operatorname{rank}_{+}}
\newcommand{\rank}{\operatorname{rank}}
\newcommand{\supp}{\operatorname{supp}}
\newcommand{\conv}{\operatorname{conv}}
\newcommand{\aff}{\operatorname{aff}}
\newcommand{\R}{\mathbb R}
\newtheorem{theorem}{Theorem}[section]
\newtheorem{lemma}[theorem]{Lemma}
\newtheorem{proposition}[theorem]{Proposition}
\newtheorem{corollary}[theorem]{Corollary}
\theoremstyle{definition}\newtheorem{definition}[theorem]{Definition}
\theoremstyle{remark}\newtheorem{remark}[theorem]{Remark}
\pagestyle{fancy}\fancyhf{}
\fancyhead[L]{\small NR-02: Cartesian-product slack matrices}
\fancyhead[R]{\small Fourth pack: partial}
\fancyfoot[C]{\thepage}
\setlength{\headheight}{14pt}
\hypersetup{pdftitle={NR-02: Certified complexity-six polygon products},pdfauthor={Sidney Holden}}
\begin{document}
\thispagestyle{empty}
\begin{center}
{\LARGE\bfseries NR-02: Cartesian-product slack matrices\par}
\vspace{9pt}
{\Large Certified complexity-six polygon products\par}
\vspace{9pt}
Sidney Holden\\{\small Center for Computational Biology, Flatiron Institute}\\{\small Simons Foundation}\\
Research note prepared with ChatGPT\\
September 13, 2026 (America/New\_York)
\end{center}
\vspace{8pt}
\noindent\fcolorbox{MidnightBlue}{MidnightBlue!4}{\parbox{0.94\textwidth}{
\textbf{Status: PARTIAL for unrestricted NR-02.} This pack gives a complete
computer-assisted argument for every pair of polygons of extension complexity
six, resolving the earlier hexagon value eleven versus twelve. It does not
prove additivity for arbitrary polytopes and does not give a counterexample.
The geometry-to-Boolean reduction is a written proof, not a proof-assistant
formalization. Independent peer review has not occurred. No novelty or
priority claim is made.}}
\begin{abstract}
Let $P,Q$ be convex polygons with $\xc(P)=\xc(Q)=6$. We prove
$\xc(P\times Q)=12$ by a finite, certificate-based support argument. Every such
polygon has between six and nine vertices. A hypothetical eleven-term
nonnegative factorization of the complete product slack matrix must satisfy
explicit row-degree, vertex-slice, edge-slice, pure-column, and coverage
conditions. These conditions give ten Boolean formulas, one for each unordered
pair of polygon sizes. Each formula has a deletion-free positive-hint LRAT
refutation checked by separate Python and C++ implementations; a further
checker recomputes unit propagation without using the hints. In total the
archived formulas contain 831,540 clauses, and the trimmed refutations contain
96,286 added clauses with 7,748,468 propagation hints. This is an exact
refutation, not a timeout or failed numerical search.

Together with the proofs preserved from the preceding packs, the result gives
additivity for every pair of polytopes with individual extension complexities
at most six. We supply an integer octagon of complexity six, its exact
six-term slack factorization, and a twelve-term factorization of its square.
We also exhibit a feasible higher-complexity support pattern that cannot admit
numerical weights: it would force incompatible antipodal facet-normal pairs.
That obstruction explains a specific limit of the present Boolean relaxation;
it does not resolve the unrestricted problem.
\end{abstract}
\textbf{Reading guide.} Sections~\ref{sec:support}--\ref{sec:certificate} are the
new lower-bound proof. Section~\ref{sec:octagon} is an exact geometric example.
Section~\ref{sec:barrier} records a genuine higher-complexity obstruction rather
than treating a support assignment as a factorization.
\clearpage
\tableofcontents
\clearpage

\section{Target and scope}
Let $P\subset\R^d$ and $Q\subset\R^e$ be full-dimensional bounded polytopes,
with complete facet--vertex slack matrices $S$ and $T$. Every irredundant facet
inequality is a row, and every vertex is a column. Define
\begin{equation}\label{eq:product}
 C_{:,(i,j)}=\begin{pmatrix}S_{:,i}\\T_{:,j}\end{pmatrix}.
\end{equation}
All factors and ranks in this note are real and nonnegative. NR-02 asks whether
\begin{equation}\label{eq:target}
 \rankp(C)=\rankp(S)+\rankp(T),
 \qquad\text{equivalently}\qquad
 \xc(P\times Q)=\xc(P)+\xc(Q).
\end{equation}
Here $\xc$ counts inequalities in a linear extension; affine equations are
free. The equivalence is the slack-factorization theorem, recalled in
\cite{twz}. The construction from two separate factorizations always gives
\begin{equation}\label{eq:upper}
 \xc(P\times Q)\le\xc(P)+\xc(Q).
\end{equation}
The repository entry rechecked for this work labels unrestricted NR-02
\emph{Partially resolved}, with a September 10, 2026 status check. It records
the proved pyramid-factor case, not a general equality \cite{repo,twz}.
This note does not alter that target or substitute arbitrary nonnegative
matrices for complete slacks.

\begin{theorem}[New standalone polygon result]\label{thm:main}
For any two convex polygons $P,Q$ with $\xc(P)=\xc(Q)=6$,
\[
 \boxed{\xc(P\times Q)=12.}
\]
The lower bound is computer-assisted, with the entire finite refutation
included in the accompanying files.
\end{theorem}

\begin{corollary}[Using the preserved previous proofs]\label{cor:small}
For any two nonempty polytopes $P,Q$ with $\xc(P),\xc(Q)\le6$,
\[
 \xc(P\times Q)=\xc(P)+\xc(Q).
\]
Points are allowed, with $\xc(\text{point})=0$.
\end{corollary}

The new standalone theorem does not depend on the previous universal
low-excess theorem. The broader corollary does use it, together with the
previous complexity-five polygon results; Section~\ref{sec:corollaries}
spells out that dependence. All source proofs are preserved in
\path{previous/NR-02_round3_proof_pack.zip}, which itself preserves the
second and first packs \cite{round3,round2}.

In particular, the unresolved interval from the preceding work,
$11\le\xc(P\times Q)\le12$ for two complexity-six hexagons, is now reduced to
$12$. The theorem also covers complexity-six heptagons, octagons, and nonagons
without imposing a particular realization. It does \emph{not} say that every
octagon or nonagon has complexity six.

\section{Why only ten polygon-size cases are required}\label{sec:finite}
We first justify the finite range without a numerical classification of
polygon realizations.

\begin{lemma}[A bounded optimal extension may be used]\label{lem:bounded}
A complete polytope slack matrix of nonnegative rank $r$ has a bounded
extension with at most $r$ facets.
\end{lemma}
\begin{proof}
Write the affine slack space as $\mathcal A=\{b-Ax:x\in\R^d\}$, so
$\mathcal A\cap\R^m_{\ge0}$ is affinely isomorphic to $P$. Boundedness gives
a strictly positive vector $\lambda$ with
$\lambda^Ts=1$ for every $s\in\mathcal A$. One way to see this is to translate
an interior point to zero, take a strictly positive dependence of the facet
normals, and normalize the resulting positive combination of right-hand sides.
The normals positively span the ambient space; summing representations of
their negatives gives a dependence with every coefficient strictly positive.

For a minimum factorization $S=WH$, every column of $W$ is nonzero. Set
\[
 E=\{h\in\R^r_{\ge0}:Wh\in\mathcal A\}.
\]
Its image is contained in the slack polytope, and the columns of $H$ lift every
vertex, so its image equals the slack polytope. Each entry of $\lambda^TW$ is
strictly positive. Consequently $\lambda^TWh=1$ bounds every coordinate of
$h$. Thus $E$ is bounded and uses at most $r$ inequalities in its affine hull.
\end{proof}

\begin{lemma}[The six-facet vertex bound]\label{lem:nine}
A polygon of extension complexity six has between six and nine vertices.
\end{lemma}
\begin{proof}
The lower bound follows from $\xc(P)\le\#\operatorname{vert}(P)$. Choose a
bounded optimal extension $E$ with six facets and dimension $h$. Since it
projects onto a polygon, $h\ge2$; boundedness gives $h\le5$.

If $h=2$, there are six vertices. If $h=3$, Euler's relation and minimum
vertex degree three give
$2f_1\ge3f_0$ and $f_0-f_1+f_2=2$, hence $f_0\le2f_2-4=8$.
If $h=5$, six facets force a simplex, with six vertices.

For $h=4$, the polar of $E$ has six vertices in four dimensions. The unique
one-dimensional space of affine dependences partitions the six indices into
$p$ positive, $q$ negative, and $t$ zero entries. Here $p,q\ge2$ and
$p+q+t=6$. A single positive entry, for example, would express that polar
vertex as a convex combination of the others, which is impossible.
The facets of the polar, and hence the vertices of $E$, correspond to the
minimal positive dependences of this one-dimensional Gale configuration.
They are exactly the opposite-sign pairs and the zero singletons. Thus
\[
 \#\operatorname{vert}(E)=pq+t\le9.
\]
For completeness, the Gale correspondence in this special case follows by
viewing a nonnegative affine function on the six polar vertices as a vector
$u\ge0$ orthogonal to the affine-dependence vector. Inclusion-minimal nonzero
supports of such vectors give maximal proper zero faces, namely facets.
Those minimal supports are precisely the pairs and singletons just counted.
The general few-facet description is also discussed in \cite{padrol}.

An affine image cannot have more vertices than its source: the inverse face
of each image vertex contains a source vertex. The upper bound nine follows
in all four dimensions.
\end{proof}

The possibilities are therefore $m,n\in\{6,7,8,9\}$. Interchanging the
factors permits $m\le n$, leaving ten cases. This is an exhaustive range for
complexity-six polygons, not a cutoff selected for computational convenience.

\section{Necessary supports of a saving factorization}\label{sec:support}
Fix $m,n\in\{6,7,8,9\}$. Number polygon edges and vertices cyclically so that
edge $f$ of $P$ contains vertices $f,f+1$, and similarly for $Q$. Indices in
one polygon are always taken modulo its size.

Suppose an eleven-term factorization of \eqref{eq:product} exists:
\[
 C=WH,\qquad W=\begin{pmatrix}A\\B\end{pmatrix},\qquad W,H\ge0.
\]
All eleven columns of $W$ may be assumed nonzero. Indeed, delete zero terms;
if the size is below eleven, split nonzero rank-one terms into positive
multiples until it is eleven. All the necessary inequalities below apply to
this padded exact factorization as well. Thus excluding eleven terms excludes
every size at most eleven.

For $f=0,\ldots,m-1$ and $g=0,\ldots,n-1$, put
\[
 K_f=\{k:A_{fk}>0\},\quad L_g=\{k:B_{gk}>0\},\quad
 \mathcal P=\bigcup_fK_f,\quad\mathcal Q=\bigcup_gL_g.
\]
The factor-index universe is $[11]=\{1,\ldots,11\}$. Define
\[
 D_P=[11]\setminus\mathcal Q,\qquad
 D_Q=[11]\setminus\mathcal P.
\]
Because every left column is nonzero, $D_P$ consists of the $P$-only columns
and $D_Q$ of the $Q$-only columns. In particular
\begin{equation}\label{eq:used}
 \mathcal P\cup\mathcal Q=[11].
\end{equation}
No numerical entry of $W$ is fixed by this notation.

\subsection{The prism bound and individual rows}
We use the established disjoint-facet inequality
\[
 \xc(R)\ge\min\{\xc(F_1),\xc(F_2)\}+2
\]
for disjoint facets of a polytope \cite{gps}. Together with the product upper
bound this gives, for every polytope $R$,
\begin{equation}\label{eq:prism}
 \xc(R\times[0,1])=\xc(R)+2.
\end{equation}
In particular, an edge times either polygon in this section has complexity
eight.

\begin{lemma}[Row supports]\label{lem:row}
Every row support is nonempty and has size at most three:
\begin{equation}\label{eq:row}
 1\le |K_f|\le3,\qquad 1\le|L_g|\le3.
\end{equation}
\end{lemma}
\begin{proof}
Every facet slack row is nonzero. On the product vertices whose $P$ coordinate
lies on edge $f$, all coefficients indexed by $K_f$ vanish. Retain the two
neighboring $P$ rows $f-1,f+1$ and all $Q$ rows. The two retained $P$ rows are
positive multiples of the two endpoint slacks of a segment. After row scaling,
this is a complete slack matrix of $[0,1]\times Q$, of nonnegative rank eight.
It uses at most $11-|K_f|$ terms. Thus $|K_f|\le3$. Interchange the factors
for the other inequality.
\end{proof}

\subsection{Pure columns and vertex restrictions}
\begin{lemma}[Pure-column caps]\label{lem:pure}
A saving factorization satisfies
\begin{equation}\label{eq:pure}
 |D_P|\le4,\qquad |D_Q|\le4.
\end{equation}
\end{lemma}
\begin{proof}
If every column were pure, the two blocks would need at least six columns
each, contradicting a size of eleven. Thus some column $k$ is mixed. Choose
a $P$ facet whose row is positive in this column and a vertex on that facet.
For this fixed $P$ vertex, coefficient row $k$ vanishes for all $Q$ vertices.
The factorization of $T$ also omits all $P$-only columns, so
$6\le11-|D_P|-1$. The mixed column is not in $D_P$, so the omissions are
distinct. This proves the first cap, and the second is symmetric.
\end{proof}

\begin{lemma}[Vertex slices]\label{lem:vertex}
For every vertex index,
\begin{equation}\label{eq:vertex}
 |K_{i-1}\cup K_i\cup D_P|\le5,
 \qquad
 |L_{j-1}\cup L_j\cup D_Q|\le5.
\end{equation}
\end{lemma}
\begin{proof}
Fix a vertex $i$ of $P$ and retain the $Q$ rows. Both incident $P$ slacks are
zero, so coefficients in $K_{i-1}\cup K_i$ vanish. The columns in $D_P$ are
absent from the retained block. All these omissions together still leave a
factorization of the rank-six matrix $T$, giving the first bound. The other
one is identical with the factors exchanged.
\end{proof}

\subsection{The strengthened edge restriction}
\begin{lemma}[Invisible and twice-killed terms]\label{lem:edge}
For every $P$ edge $f$,
\begin{equation}\label{eq:edge}
 \left|K_f\ \cup\ (K_{f-1}\cap K_{f+1})\ \cup\
 \bigl(D_P\setminus(K_{f-1}\cup K_{f+1})\bigr)\right|\le3.
\end{equation}
The analogous assertion holds for the $L$ supports and $D_Q$.
\end{lemma}
\begin{proof}
Use the edge test from Lemma~\ref{lem:row}, retaining only the neighboring
$P$ rows and all $Q$ rows. Terms in $K_f$ vanish throughout the test. A term
in both neighboring supports has zero coefficient at each of the two
endpoints: one neighbor forces it to zero at one endpoint and the other
neighbor does so at the other. Finally, a $P$-only column absent from both
neighbors is zero on every retained row, whether or not its coefficient
vanishes. The union of these three sets can therefore be omitted. The
remaining matrix has nonnegative rank eight, so at most three terms are
omitted. Counting a union is essential; overlapping reasons are not added
as separate losses.
\end{proof}

\subsection{Positive-entry coverage and mixed facet unions}
For a product vertex $(i,j)$ define its forced-zero support union
\[
 Z_{ij}=K_{i-1}\cup K_i\cup L_{j-1}\cup L_j.
\]
Every positive entry of the product must be contributed by some index outside
this union. Consequently,
\begin{align}
 K_f\setminus Z_{ij}&\ne\varnothing
       &&\text{if }i\notin\{f,f+1\},\label{eq:coverageP}\\
 L_g\setminus Z_{ij}&\ne\varnothing
       &&\text{if }j\notin\{g,g+1\}.\label{eq:coverageQ}
\end{align}
These assertions are necessary even though a coefficient outside $Z_{ij}$
need not actually be positive. Allowing all such coefficients enlarges the
feasible support family and is therefore safe for an impossibility proof.

\begin{lemma}[Cross-union restriction]\label{lem:cross}
For every pair $K_f,L_g$, their union contains no third facet-row support:
\begin{equation}\label{eq:cross}
 K_h\not\subseteq K_f\cup L_g\ (h\ne f),\qquad
 L_t\not\subseteq K_f\cup L_g\ (t\ne g).
\end{equation}
\end{lemma}
\begin{proof}
At a point of $\operatorname{relint}(F_f)\times\operatorname{relint}(G_g)$,
exactly the named two product facet slacks vanish. Express the point as a
convex combination of product vertices and use the same combination of
coefficient columns as a nonnegative lift. Every coefficient indexed by
$K_f\cup L_g$ is zero. A third support contained in the union would force a
third slack to vanish, a contradiction.
\end{proof}

This last condition is redundant with full positive-entry coverage, but
including it provides useful short deductions in the finite proof. Its use
does not add a geometric hypothesis.

\section{The exact Boolean formulas}\label{sec:encoding}
Let $x_{fk}$ be a Boolean variable indicating membership of factor index $k$
in combined left row $f$, with the $m$ rows of $P$ followed by the $n$ rows
of $Q$. The primary variables number $11(m+n)$. The function
\texttt{build(m,n)} in \texttt{code/support\_cnf.py} encodes exactly the
necessary conditions of Section~\ref{sec:support}, plus column-permutation
symmetry. Its default arguments are $a=b=6$, $r=11$, and no fixed row degree.
The optional generalized arguments are used only in the explicitly labeled
controls and higher-complexity experiments.

\subsection{Gates, bounds, and witnesses}
For a Boolean output $y=\bigwedge_{s=1}^t\ell_s$, where the $\ell_s$ are signed
literals, the encoder uses
\[
 (\neg y\lor\ell_s)\quad(s=1,\ldots,t),\qquad
 (y\lor\neg\ell_1\lor\cdots\lor\neg\ell_t).
\]
OR gates use the negation of an AND of the negated inputs. Empty AND inputs
mean true. All gate clauses are exact equivalences.

The constraint that at most $k$ of $\ell_1,\ldots,\ell_r$ hold is encoded
by the clause $\bigvee_{s\in I}\neg\ell_s$ for every $(k+1)$-element subset
$I$. At the proof parameters, $r=11$ and the relevant caps are three, four,
and five. These are explicit cardinality clauses, not a floating-point
integer-programming relaxation.

For a positive entry, a witness $y_k$ implies its row membership and the
absence of the four forced-zero row memberships at index $k$. The encoder
requires $\bigvee_k y_k$. The reverse implications into a witness are not
needed: whenever the underlying set condition holds, choose a valid witness.
The cross-union exclusions use the same construction with two excluded rows.

\subsection{Column symmetry and equisatisfiability}
Factor columns may be permuted freely. We sort their complete binary row
incidence vectors in lexicographically nonincreasing order, reading row zero
first. Equal support vectors are permitted.

For adjacent columns with current bits $a,b$ and a Boolean $e$ saying that
all previous bits agreed, the order condition is
$\neg e\lor a\lor\neg b$. A new exact gate records
$e'=e\wedge(a=b)$; the initial prefix is true. The supplied six-clause
implementation of this step is checked by exhaustive small comparator tests.

Every genuine factorization satisfying the hypotheses can have its columns
sorted, producing a primary assignment satisfying these constraints. Set all
gate variables to their defining values and choose valid positive witnesses;
this extends it to a satisfying assignment of the full CNF. Conversely, the
gate clauses and witness implications recover every stated set condition
from any satisfying assignment. Hence the encoding neither loses a possible
factorization through its auxiliary variables nor presumes numerical weights
for a merely feasible support assignment.

The method needs only the forward implication from a factorization to a
satisfying Boolean assignment. It never asserts the converse from Boolean
supports to a nonnegative factorization.

\section{The finite refutations and their verification}\label{sec:certificate}
\begin{proposition}[Certified finite support exclusion]\label{prop:finite}
For each $6\le m\le n\le9$, the default formula generated by
\texttt{build(m,n)} is unsatisfiable.
\end{proposition}
\begin{proof}[Computer-assisted proof]
For each pair, \texttt{proofs/polygon\_m\_n/instance.cnf} is the exact generated
formula, and \texttt{proof.lrat} is a complete deletion-free positive-hint
refutation. The following are the checked certificate sizes. Each row ends
with a checked empty clause.
\begin{center}\small
\begin{tabular}{@{}ccrrrr@{}}\toprule
$m$&$n$&Variables&Source clauses&Added clauses&Unit hints\\\midrule
6&6&7,940&45,643&5,439&346,981\\
6&7&10,117&55,510&9,614&703,887\\
6&8&12,558&66,457&7,093&428,745\\
6&9&15,263&78,484&7,019&533,361\\
7&7&12,800&67,425&15,447&1,048,486\\
7&8&15,791&80,600&9,987&757,969\\
7&9&19,090&95,035&13,920&1,111,358\\
8&8&19,376&96,183&8,144&826,076\\
8&9&23,313&113,206&7,462&851,199\\
9&9&27,932&132,997&12,161&1,140,406\\\midrule
\multicolumn{3}{r}{Totals}&831,540&96,286&7,748,468\\\bottomrule
\end{tabular}
\end{center}
The Python and C++ direct-hint checkers independently accepted every step.
A further C++ checker discarded the hints and recomputed reverse unit
propagation for all added clauses. The separate occurrence-list checker gives
another propagation implementation. The verifier also regenerates every CNF
and compares it byte-for-byte, rather than trusting an unrelated Boolean
file with the same title.

Here is the precise proof rule. To check an added clause $D$, temporarily
assume all literals of $D$ false. Each positive hint names an earlier source
or proved clause. After the current assignments, that clause must be unit
or empty. In the unit case, assign its last literal true; in the empty case,
a contradiction has been reached. A step without a contradiction is rejected.
Every hinted clause must already exist, so forward references and new axioms
are forbidden. The proof must eventually add the empty clause.

This rule is sound: any model of the preceding clauses that falsified $D$
would have to satisfy each unit implication in sequence, and hence would
satisfy a contradiction. Therefore the preceding formula entails $D$.
Induction over the added clauses proves that a checked empty clause implies
unsatisfiability of the original formula. The archived proof contains only
this positive-hint fragment of LRAT \cite{lrat}; no RAT step or deletion is
needed for the final verification.
\end{proof}

\begin{proof}[Proof of Theorem~\ref{thm:main}]
The product upper bound gives twelve. If the product had nonnegative rank at
most eleven, pad an exact factorization to eleven nonzero terms as described
in Section~\ref{sec:support}. Lemma~\ref{lem:nine} puts its polygon sizes in
one of the ten cases, after exchanging the factors if necessary.
Every support constraint follows from that factorization. Sorting its columns
and assigning the auxiliary variables would satisfy the corresponding CNF.
Proposition~\ref{prop:finite} contradicts this. Thus the nonnegative rank is
at least twelve, giving equality.
\end{proof}

\subsection{What is and is not trusted}
The discovery solver and the proof-trimming producer are not needed to verify
the archived conclusion. Their outputs are checked using the original CNF
and explicit clauses. Some preliminary discovery traces did not pass our
restrictive RUP-only checker; they were discarded, not treated as proofs.
Only accepted traces were converted, trimmed, and checked again.

The two direct-hint implementations were written separately in this work,
but they are not independent human reviews or formally verified programs.
The geometric implication from a slack factorization to the Boolean formula
is the written proof in Section~\ref{sec:support}, not a Lean theorem. A
correct refutation of a wrongly encoded problem would not prove the desired
mathematical claim. For this reason the package includes the encoder,
independent set-condition audits, positive controls, negative checker tests,
and a detailed proof-audit checklist.

Unlike the earlier bounded searches, Proposition~\ref{prop:finite} is an
exhaustive deduction: every possible Boolean assignment is excluded by a
checked refutation. It is not an inference from failure to find an assignment
or from a small residual in numerical optimization.

\section{Consequences and their dependencies}\label{sec:corollaries}
\begin{corollary}[All hexagon pairs]
For every two convex hexagons,
$\xc(P\times Q)=\xc(P)+\xc(Q)$.
\end{corollary}
\begin{proof}
A hexagon has complexity five or six. If both values are six, apply
Theorem~\ref{thm:main}. The cases involving a value five were proved in the
second pack, Corollaries 7.3 and 7.6, whose complete proofs are preserved
\cite{round2}.
\end{proof}

\begin{proof}[Proof of Corollary~\ref{cor:small}]
A point or segment is additive with every factor. If either factor has
dimension at least three, its complexity at most six is at most its dimension
plus three. The universal low-excess theorem from the third pack applies
\cite{round3}.

It remains to consider polygons of complexities three through six. A triangle
is a pyramid and is universally additive \cite{twz}. A complexity-four polygon
is a quadrilateral. Two disjoint edge-product facets and the prism bound give
universal additivity for a quadrilateral as well.

A complexity-five polygon has five or six vertices: a bounded five-facet
extension has at most six vertices, by the dimension-two, three, and four
counts used below. Thus it is a pentagon or a complexity-five hexagon. The
second pack proves pentagon additivity against any factor of complexity at
most seven, and complexity-five-hexagon additivity against any factor of
complexity at most eight \cite{round2}. These cover every remaining pair with
at least one value five. Finally the pair of values six is
Theorem~\ref{thm:main}.
\end{proof}

The small-complexity corollary excludes a counterexample with both individual
complexities at most six. It does not make complexity-six polygons universally
additive against arbitrary other factors. Nor does pairwise additivity justify
arbitrary powers: after multiplying two complexity-six factors the intermediate
complexity is twelve, outside the hypothesis. The previous arbitrary-pentagon-
product theorem remains a separate result \cite{round3}.

\section{An integer octagon and its exact square}\label{sec:octagon}
The following example gives an eight-vertex member of the new theorem without
assuming a general polygon classification. Let $P$ have the cyclic vertices
\begin{equation}\label{eq:octvertices}
\begin{split}
&(-64,-64),\ (-32,-57),\ (32,-41),\ (64,-32),\\
&(64,32),\ (32,41),\ (-32,57),\ (-64,64).
\end{split}
\end{equation}
In coordinates $(v,w)$ its eight facet slacks are the rows of
\[
 F\begin{pmatrix}1\\v\\w\end{pmatrix},\qquad
 F=\begin{pmatrix}
1600&-7&32\\196&-1&4\\1600&-9&32\\64&-1&0\\
1600&-9&-32\\196&-1&-4\\1600&-7&-32\\64&1&0
\end{pmatrix}.
\]
At the listed vertices every slack is nonnegative and exactly the two incident
facet slacks are zero. The verifier also enumerates all feasible pairwise
facet intersections, obtaining exactly these eight vertices.

\subsection{A six-facet extension}
In coordinates $(t,v,w)$, take the six nonnegative affine functions
\begin{align*}
 h_0&=1+t,& h_1&=1-t,\\
 h_2&=48-16t+v,& h_3&=48-16t-v,\\
 h_4&=194+2t-v+4w,& h_5&=194+2t-v-4w.
\end{align*}
Their simultaneous nonnegativity defines a bounded three-dimensional
polytope $E$. Indeed $-1\le t\le1$, $|v|\le48-16t$, and
$4|w|\le194+2t-v$. The last right-hand side is strictly positive throughout
the first two bounds. Successive endpoint choices give precisely eight
vertices; direct enumeration of all triples of facet equations confirms this.
Each of the six inequalities defines a genuine facet.

Project by $(t,v,w)\mapsto(v,w)$. The lifts of \eqref{eq:octvertices} use
$t=-1$ at $v=\pm64$ and $t=1$ at $v=\pm32$. Thus the image is exactly $P$,
showing $\xc(P)\le6$.

For the reverse inequality, any bounded extension with at most five facets
has dimension at most four. In dimension two it has at most five vertices;
in dimension three the Euler bound gives at most six; in dimension four it
is a simplex with five vertices. None can project to eight vertices.
Lemma~\ref{lem:bounded} therefore yields
\[
 \boxed{\xc(P)=6.}
\]

\subsection{Exact slack factors}
Let $H$ consist of the six extension slacks evaluated at the eight lifts, in
the order of \eqref{eq:octvertices}. The polygon slack matrix is $S=WH$ with
\[
 W=\begin{pmatrix}
0&0&1&0&8&0\\
0&2&0&0&1&0\\
0&0&0&1&8&0\\
16&0&0&1&0&0\\
0&0&0&1&0&8\\
0&2&0&0&0&1\\
0&0&1&0&0&8\\
16&0&1&0&0&0
\end{pmatrix}.
\]
All entries of $W,H,S$ are integers. The exact computed ranks are
$\rank(S)=3$, $\rank(W)=5$, and $\rank(H)=4$; they are not substitutes for the
nonnegative-rank lower bound just proved.

The complete square-product slack matrix $C$ has size $16\times64$, with
\[
 C_{:,(i,j)}=\begin{pmatrix}S_{:,i}\\S_{:,j}\end{pmatrix},\qquad
 W_{\square}=\operatorname{diag}(W,W),\qquad
 (H_{\square})_{:,(i,j)}=\begin{pmatrix}H_{:,i}\\H_{:,j}\end{pmatrix}.
\]
The data file gives every entry and verifies $W_{\square}H_{\square}=C$
exactly. The ordinary rank is five. Theorem~\ref{thm:main} gives the matching
nonnegative-rank lower bound, so
\[
 \boxed{\rankp(C)=\xc(P\times P)=12.}
\]
This is a solved exact instance beyond the original hexagon examples.

\section{A higher-complexity support pattern is not a counterexample}\label{sec:barrier}
The same necessary conditions can be generated with assumed individual
complexities $a,b$ and a saving size $r=a+b-1$. The row caps become
$a-3$ and $b-3$, vertex omission caps become $a-1$ and $b-1$, and the pure-column
caps become $a-2$ and $b-2$. The strengthened edge formula is unchanged apart
from those corresponding caps.

For a putative hexagon of complexity six and an octagon of complexity seven,
with $r=12$, this generalized support system is satisfiable. The saved Boolean
model has each hexagon row supported on its own pair of columns
$\{2g,2g+1\}$, $g=0,\ldots,5$. The two octagon supports of those columns are
as follows, using zero-based row labels:
\begin{center}
\begin{tabular}{@{}ccc@{}}\toprule
Hexagon label $g$&Octagon support of $2g$&Octagon support of $2g+1$\\\midrule
0&$\{0,1,2\}$&$\{4,5,6\}$\\
1&$\{0,6,7\}$&$\{2,3,4\}$\\
2&$\{1,5\}$&$\{3,7\}$\\
3&$\{0,1,2\}$&$\{4,5,6\}$\\
4&$\{1,5\}$&$\{3,7\}$\\
5&$\{0,6,7\}$&$\{2,3,4\}$\\\bottomrule
\end{tabular}
\end{center}
The package verifies every generalized support condition for this pattern
using integer bit sets. Nevertheless the pattern cannot be assigned positive
left-factor weights for any complete hexagon--octagon slack product.

\begin{proposition}[Nonalternating forced antipodes]\label{prop:barrier}
There is no product slack factorization with exactly the positive left-column
supports in the preceding table and the indicated hexagon pairs.
\end{proposition}
\begin{proof}
Normalize each left column so its single positive hexagon-block entry equals
one. For every hexagon facet $g$, the two coefficient rows then satisfy
\[
 H_{2g,(i,j)}+H_{2g+1,(i,j)}=S_{g,i}.
\]
For every octagon vertex $j$, the zero octagon rows kill exactly one member
of each pair, as a direct reading of the table shows. Thus the other member
is forced to equal $S_{g,i}$; the split is not a free numerical variable.

At octagon vertex $j=2$, row zero receives contributions only from columns
$2$ and $10$. Therefore, for some $\alpha,\beta>0$,
\[
 T_{0,2}=\alpha S_{1,i}+\beta S_{5,i}\quad\text{for every }i.
\]
At octagon vertex $j=0$, row one similarly gives
\[
 T_{1,0}=\gamma S_{2,i}+\delta S_{4,i}\quad\text{for every }i,
 \qquad\gamma,\delta>0.
\]
Both left-hand sides are strictly positive constants. Since the hexagon
vertices affinely span the plane, these equalities extend to affine identities.
Writing its facet slacks as $s_f(x)=b_f-n_f^Tx$ yields
\[
 \alpha n_1+\beta n_5=0,\qquad
 \gamma n_2+\delta n_4=0.
\]
Thus the two facet-normal pairs $(1,5)$ and $(2,4)$ must be antipodal.
Outward normal rays of an irredundant convex polygon are distinct and occur
in strict cyclic facet order. Two distinct antipodal pairs must alternate
in that order. These pairs do not: their order is $1,2,4,5$. More explicitly,
choose angles with $\theta_1<\theta_2<\theta_4<\theta_5<\theta_1+2\pi$.
Antipodality would require $\theta_5=\theta_1+\pi$ and
$\theta_4=\theta_2+\pi>\theta_5$, a contradiction.
\end{proof}

The proposition excludes this particular satisfying support assignment, not
all size-twelve candidates at assumed ranks $(6,7)$. A second support-only
assignment at assumed ranks $(7,7)$ is also recorded, without a numerical
realization certificate. Neither is an NR-02 counterexample.

This supplies a concrete boundary of the method: the lower-complexity
refutations close the support problem completely, but support feasibility at
the next complexities may still violate geometric identities. A general proof
or counterexample needs control of these numerical and geometric conditions,
not just a larger unweighted cover search.

\section{Reproducibility and verification scope}\label{sec:verification}
The full command from the extracted archive is
\begin{verbatim}
python -m pip install -r requirements.txt
python code/verify_all.py --require-cpp
\end{verbatim}
The LRAT checker itself requires only Python's standard library. SymPy is
used for exact geometric data, and NumPy is required by the preserved older
packs. No SAT solver or network access is used in verification.

The optional deeper propagation check is
\begin{verbatim}
python code/verify_all.py --require-cpp --deep-rup
\end{verbatim}
It additionally uses the independently written occurrence-list RUP checker
on every refutation, instead of trusting either the hints or the watched-
literal implementation. Required checks fail rather than being silently
reported as passed. Run without Python's optimization flag because the older
packs use exact assertions.

The accompanying tests include the following distinct layers:
\begin{center}\small
\begin{tabular}{@{}p{0.30\textwidth}p{0.63\textwidth}@{}}\toprule
Layer&Exact scope\\\midrule
Finite lower bound&All ten CNF formulas and their complete checked refutations;
96,286 added clauses ending in ten empty clauses.\\
Encoding truth tables&254 AND-gate assignments; 16 prefix-gate assignments;
6,144 cardinality assignments;
5,460 comparator pairs.\\
Set/CNF comparison&600 complete primary assignments and 5,400 individual
constraint-group comparisons across all ten formula sizes.\\
Positive encoding controls&A genuine twelve-term octagon-square factorization
satisfies the nonsaving control formula; the higher-rank support assignment
satisfies its generalized formula.\\
Checker controls&20 valid toy refutations accepted and 23 false, malformed or
incomplete refutations rejected by the available implementations; 150 toy
assignments checked exhaustively.\\
Exact geometric data&All octagon and extension vertices/facets, $WH=S$,
$W_{\square}H_{\square}=C$, and exact ordinary ranks.\\
Earlier work&The third pack's verifier, including the second pack and all first-
pack checks, is rerun from the unchanged preserved archive.\\\bottomrule
\end{tabular}
\end{center}
The finite sample checks in the encoding audit do not replace the universal
encoding argument. Likewise, matrix multiplication checks establish the upper
factorizations, not their nonnegative-rank lower bounds. The latter come from
the geometric vertex count and the finite support refutations.

\texttt{SHA256SUMS.txt} detects altered package files. It is not a proof of
mathematical truth. The archive includes source and data, not compiled
executables or font files. The final PDF source is rebuildable with two
\texttt{pdflatex} passes.

\section{What remains unresolved}\label{sec:remaining}
The general equality \eqref{eq:target} is still not established. The new result
settles the previously identified complexity-six hexagon pair and the broader
bounded-complexity class, but it gives no reduction of arbitrary polytopes to
that class.

In particular, the earlier pentagon-factor residual regime remains:
\[
 a=\xc(P)\ge\max\{8,\dim P+6\},\qquad
 \xc(P\times Q)=a+4
\]
for a pentagon $Q$, with the paired form and higher-dimensional slice fibers
from the second pack. The complexity-five-hexagon residual regime is also
unchanged. The general geometry of dense factor rows and high-dimensional
optimal extensions is not eliminated here.

The higher-complexity models in Section~\ref{sec:barrier} are not weighted
factorizations and therefore are not counterexamples. Conversely, the
nonalternating-antipodes argument proves only that one of those support
patterns is unrealizable; it is not a universal lower bound at ranks $(6,7)$.

A full negative resolution still requires complete slack matrices, an exact
saving factorization, and rigorous individual nonnegative-rank lower bounds.
A full positive resolution still requires a lower-bound mechanism covering
all factor complexities and dimensions. Neither is supplied in this pack.
The appropriate unrestricted status remains a partial contribution, with a
complete computer-assisted proof of the stated restricted theorem. The
repository's distinction between partial contributions, claimed resolutions,
and proof-assistant verification should be maintained \cite{guidelines}.

\clearpage
\phantomsection
\addcontentsline{toc}{section}{References}
\begin{thebibliography}{9}
\bibitem{repo} OpenProblemsInNLA. NR-02, ``Additivity for Cartesian-product slack
matrices.'' Entry rechecked September 13, 2026; recorded status check
September 10, 2026.
\url{https://github.com/ajt60gaibb/OpenProblemsInNLA/blob/main/nonnegative-and-positive-factorizations/NR-02/README.md}.
\bibitem{twz} H. R. Tiwary, S. Weltge, and R. Zenklusen. Extension complexities
of Cartesian products involving a pyramid. \emph{Information Processing
Letters} 128 (2017), 11--13. The slack-factorization theorem is recalled in
Section 2; the product theorem concerns a pyramid factor.
\url{https://arxiv.org/abs/1702.01959}.
\bibitem{gps} F. Grande, A. Padrol, and R. Sanyal. Extension complexity and
realization spaces of hypersimplices. arXiv:1601.02416v2 (2017). Lemma 2.3 and
Corollary 2.4 supply the disjoint-facet and prism bounds.
\url{https://arxiv.org/html/1601.02416v2}.
\bibitem{padrol} A. Padrol. Extension complexity of polytopes with few vertices
or facets. arXiv:1602.06894v2 (2016). The few-facet discussion provides the
relevant one-dimensional Gale description.
\url{https://arxiv.org/html/1602.06894v2}.
\bibitem{lrat} L. Cruz-Filipe, M. J. H. Heule, W. A. Hunt Jr., M. Kaufmann,
and P. Schneider-Kamp. Efficient certified RAT verification.
arXiv:1612.02353 (2016/2017).
\url{https://arxiv.org/abs/1612.02353}.
The present checkers use only positive unit-propagation hints; the formally
verified checkers described in that paper were not run in this work.
\bibitem{round3} \emph{NR-02: Universal low-excess factors and additive face
envelopes}. Third research note prepared in this conversation, September 13,
2026. Preserved unchanged in \path{previous/NR-02_round3_proof_pack.zip}.
This is prior author-side work, not an independent publication.
\bibitem{round2} \emph{NR-02: Stronger pentagon bounds and sparse-factor
obstructions}. Second research note prepared in this conversation, September
13, 2026. Preserved within the third pack, together with the first pack.
\bibitem{guidelines} OpenProblemsInNLA. Contribution and reporting guidance.
\url{https://github.com/ajt60gaibb/OpenProblemsInNLA/blob/main/CONTRIBUTING.md}.
\end{thebibliography}
\end{document}
